
\section{An explicit nonzero-data unforced comparison and a sharp obstruction to small corrections}
\label{sec:xxxiv-explicit-comparison}

The central remaining step is not to produce an arbitrary
unforced flow with nonzero data---such flows are classical---but to
prove that its correction to the \emph{actual} forced singular
flow is smaller than the source at the collapsing core.
We construct here a concrete global whole-space unforced comparison,
with independently certifiable critical Fourier smallness and
a complete zero-force PDE identity. We then prove that its correction
necessarily \emph{violates} the proposed
$o(\tau^{-4A/3})$ $H^2$ bound. This is an adversarial
test of the conjectured cancellation mechanism, not a proof of
finite-time unforced Navier--Stokes blowup.

\begin{theorem}[Explicit smooth nonzero-data unforced whole-space flow]
\label{thm:xxxiv-gaussian}
Let $\nu>0$, $\rho^2=x^2+y^2+z^2$ and prescribe
\begin{equation}\label{eq:xxxiv-datum}
\boxed{
v_0(x,y,z)=\frac{\nu}{100\pi}
(-2y e^{-\rho^2},\,2x e^{-\rho^2},\,0).
}
\end{equation}
There exists a global smooth divergence-free, finite-energy
solution $v$ of the exact \emph{unforced} Navier--Stokes system
on $\mathbb R^3$ with datum $v_0$. In particular, with
\[
C_{\rm F}=(2\pi)^{-3/2},\qquad
\|a\|_{\mathcal X^j}=\int_{\mathbb R^3}
|\xi|^j|\widehat a(\xi)|\,d\xi
\]
and the unitary Fourier transform, the explicit data satisfy
\[
\boxed{
C_{\rm F}\frac{\|v_0\|_{\mathcal X^{-1}}}{\nu}
=\frac1{400},\quad
\|v_0\|_2=\frac{\nu}{100(2\pi)^{1/4}}.
}
\]
One may construct $v$ by a globally convergent Fourier mild
iteration with contraction factor at most $1/25$.
For all $t\ge0$, the comparison also satisfies
\begin{equation}\label{eq:xxxiv-uniform-bound}
\boxed{\quad
\|v(t)\|_\infty\le
C_v:=\frac{\nu}{100\sqrt\pi}e^{1/50}.\quad}
\end{equation}
Its pressure may be fixed by
$p_v=\sum_{i,j}R_iR_j(v_iv_j)$, where $R_i$ are the
standard Riesz transforms; then
\[
\partial_tv+(v\cdot\nabla)v-\nu\Delta v+\nabla p_v=0
\]
holds identically for every $t>0$ and by smooth continuity
at $t=0$.
\end{theorem}
\begin{proof}
Let $\psi=e^{-\rho^2}$. The datum is
$(\nu/(100\pi))(\partial_y\psi,-\partial_x\psi,0)$,
so it is smooth, solenoidal, nonzero and Schwartz.
The unitary transform is
$\widehat\psi(\xi)=2^{-3/2}e^{-|\xi|^2/4}$. Thus
$|\widehat{\nabla\times(\psi e_z)}|
=2^{-3/2}\sqrt{\xi_x^2+\xi_y^2}\,
 e^{-|\xi|^2/4}$, and exact spherical integration gives
\begin{equation}\label{eq:xxxiv-gaussian-norms}
\|\nabla\times(\psi e_z)\|_{\mathcal X^{-1}}
=\frac{\pi^{5/2}}{\sqrt2},\qquad
\|\nabla\times(\psi e_z)\|_{\mathcal X^0}
=2\sqrt2\pi^2,\qquad
\|\nabla\times(\psi e_z)\|_2^2
=2(\pi/2)^{3/2}.
\end{equation}
The critical smallness ratio is consequently $1/400$.

For completeness, define the time norm
\[
\|u\|_{\mathscr B}
=\sup_{t\ge0}\|u(t)\|_{\mathcal X^{-1}}
+\nu\int_0^\infty\|u(t)\|_{\mathcal X^1}\,dt.
\]
With $H(t)=e^{\nu t\Delta}v_0$, direct integration yields
$\|H\|_{\mathscr B}\le2X_0$, where
$X_0=\|v_0\|_{\mathcal X^{-1}}$.
Let
\[
\mathcal B(u,v)(t)=-
\int_0^t e^{\nu(t-s)\Delta}
\mathbb P\nabla\cdot(u\otimes v)(s)\,ds.
\]
The Fourier multiplier and $\|\mathbb P_\xi\|_{\rm op}\le1$
give for both the supremum $\mathcal X^{-1}$ norm and the
time-integrated $\nu\mathcal X^1$ norm a bound by
$C_{\rm F}\int\|u\|_{\mathcal X^0}
\|v\|_{\mathcal X^0}\,dt$.
Cauchy--Schwarz and
$\|u\|_{\mathcal X^0}^2
\le\|u\|_{\mathcal X^{-1}}\|u\|_{\mathcal X^1}$
therefore give
\[
\|\mathcal B(u,v)\|_{\mathscr B}
\le\frac{2C_{\rm F}}{\nu}
\|u\|_{\mathscr B}\|v\|_{\mathscr B}.
\]
The map $u\mapsto H+\mathcal B(u,u)$ sends the closed
ball $\|u\|_{\mathscr B}\le4X_0$ into itself, since
$2X_0+32(C_{\rm F}X_0/\nu)X_0<4X_0$,
and has Lipschitz constant at most
$16C_{\rm F}X_0/\nu=1/25$ there.
Banach's theorem gives a global divergence-free mild
solution, satisfying the projected PDE. Parabolic smoothing
and the Schwartz initial datum give a classical solution;
the pressure formula reconstructs the exact physical zero-force
equation. Standard energy estimates yield
$\|v(t)\|_2\le\|v_0\|_2$.

The Fourier $\mathcal X^0$ inequality is
\[
\frac{d}{dt}\|v(t)\|_{\mathcal X^0}
\le2C_{\rm F}
\|v(t)\|_{\mathcal X^0}\|v(t)\|_{\mathcal X^1}.
\]
Since $\nu\int_0^\infty\|v\|_{\mathcal X^1}
\le4X_0$, Gronwall implies
$\|v(t)\|_{\mathcal X^0}
\le\|v_0\|_{\mathcal X^0}
 e^{8C_{\rm F}X_0/\nu}
=\|v_0\|_{\mathcal X^0}e^{1/50}$.
Finally Fourier inversion and
$C_{\rm F}\|v_0\|_{\mathcal X^0}
=\nu/(100\sqrt\pi)$ give
\eqref{eq:xxxiv-uniform-bound}.
This boundedness, combined with standard Sobolev energy
estimates, excludes any finite-time loss of smoothness.
\end{proof}

\begin{theorem}[Exact correction equation and nonintegrable $H^2$ cancellation]
\label{thm:xxxiv-h2-obstruction}
Let $u_B$ be the actual compactly supported forced
velocity of OpenAI \cite{openai2026}, Theorem~1.1,
with viscosity $\nu$ and complete physical force $f_B$.
Put $w=v-u_B$, where $v$ is the globally smooth
\emph{unforced} solution of
Theorem~\ref{thm:xxxiv-gaussian}. Then
\[
w(0)=v_0\ne0,\qquad
\partial_tw-\nu\Delta w+
\mathbb P\nabla\cdot
(u_B\otimes w+w\otimes u_B+w\otimes w)
=-\mathbb P f_B.
\]
Thus $u_B+w=v$ satisfies the \emph{full zero-force PDE
identity with reconstructed pressure}; no projected-force
or cutoff term has been omitted.

Let $A=1/2+h$, $e_0>0$ and $x_\tau$ be the source's
actual Section~10 singular trajectory,
and write
\[
B=\sup_{0\le t<1}\|u_B(t)\|_2,\qquad
M=B+\frac{\nu}{100(2\pi)^{1/4}},\qquad
C_2=\frac4{\pi\sqrt2\,3^{7/8}}.
\]
For all sufficiently small $\tau>0$,
\begin{equation}\label{eq:xxxiv-point-cancel}
|w_\theta(x_\tau,1-\tau)|
\ge\frac{e_0}{4}\tau^{-A},
\end{equation}
and the correction necessarily obeys
\begin{equation}\label{eq:xxxiv-h2-nogo}
\boxed{\quad
\|\Delta w(1-\tau)\|_2
\ge c_{\nu,B,e_0}\tau^{-4A/3},\qquad
c_{\nu,B,e_0}:=
\left(\frac{e_0}{4C_2 M^{1/4}}\right)^{4/3}>0.
\quad}
\end{equation}
Consequently it is \emph{impossible} for this exact
nonzero-data, zero-force corrector to satisfy
$\|\Delta w(1-\tau)\|_2=o(\tau^{-4A/3})$,
and indeed
\begin{equation}\label{eq:xxxiv-L2time-nogo}
\boxed{\quad
\int_{t_0}^{1}\|\Delta w(t)\|_2^2\,dt=\infty
\qquad\text{for every }t_0<1.\quad}
\end{equation}
\end{theorem}
\begin{proof}
Subtract the source's fully forced equation from the exact
zero-force physical equation for $v$; the projected displayed
difference equation follows identically.
For sufficiently small $\tau$, the source (10.21) yields
$|u_{B,\theta}(x_\tau,1-\tau)|
\ge(e_0/2)\tau^{-A}$. On the other hand
$|v_\theta|\le\|v\|_\infty\le C_v$.
Choose $\tau$ small enough that
$C_v\le(e_0/4)\tau^{-A}$ and apply the reverse
triangle inequality to obtain \eqref{eq:xxxiv-point-cancel}.
By the energy inequality and the source's uniform $L^2$
bound, $\|w(t)\|_2\le M$.
The explicit three-dimensional Fourier interpolation inequality
of Lemma~\ref{lem:xxxiii-agn} then gives
\[
|w(x_\tau,1-\tau)|
\le C_2 M^{1/4}\|\Delta w(1-\tau)\|_2^{3/4}.
\]
This proves \eqref{eq:xxxiv-h2-nogo}.
Since $8A/3>4/3>1$, its square is not integrable
in a neighborhood of $\tau=0$, proving
\eqref{eq:xxxiv-L2time-nogo}.
\end{proof}

\paragraph{Scientific interpretation and exact unresolved target.}
A smooth, numerically specified nonzero datum and a global
exactly unforced whole-space solution have been constructed;
the complete physical zero-force identity is therefore
established \emph{for that completed flow}. But that solution
is regular and its correction exactly cancels the forced
singularity. The required little-$o$ $H^2$ bound is
\emph{false for this concrete correction}. Proving it for
a different exact unforced correction while preserving
the published source core would amount to a finite-time
singularity of smooth-data unforced Navier--Stokes---the
unresolved Millennium regularity problem.
The Fourier smallness constants, contraction factor,
uniform amplitude, and $H^2$ time-divergence have been
independently recomputed; finite Lean scalar dependencies
are separate from a full PDE formalization.


\begin{proposition}[A genuine uniform two-sided $H^2$ estimate for the explicit correction]
\label{prop:xxxiv-two-sided}
For the exact nonzero-data unforced comparison in
Theorem~\ref{thm:xxxiv-gaussian}, the standard ordered-derivative
$H^2$ norm satisfies for every $t\ge0$
\[
\boxed{\quad
\|v(t)\|_{H^2}\le
C_{H,\nu}:=
\frac{\nu\sqrt{82}}{100\pi}
\left(\frac\pi2\right)^{3/4}e^{3/10}.\quad}
\]
Consequently for its actual source correction $w=v-u_B$,
\begin{equation}\label{eq:xxxiv-h2-two-sided}
\boxed{\quad
\max\{0,\|\Delta u_B(t)\|_2-C_{H,\nu}\}
\le\|\Delta w(t)\|_2
\le\|\Delta u_B(t)\|_2+C_{H,\nu}.\quad}
\end{equation}
Moreover, as $t\uparrow1$,
\begin{equation}\label{eq:xxxiv-h2-ratio}
\boxed{\quad
\frac{\|\Delta w(t)\|_2}{\|\Delta u_B(t)\|_2}
\longrightarrow1.\quad}
\end{equation}
This is an actual $H^2$ upper bound for an exact nonzero-data
unforced correction, but it is \emph{not} the desired
little-$o$ upper bound: it proves that this comparison
inherits essentially all of the forced source's $H^2$
divergence.
\end{proposition}
\begin{proof}
Use the equivalent squared $H^2$ norm
$\|v\|_2^2+\sum_i\|\partial_iv\|_2^2+
\sum_{i,j}\|\partial_i\partial_jv\|_2^2$.
In each derivative energy identity, the leading
$v\cdot\nabla\partial^\alpha v$ term vanishes because
$\nabla\cdot v=0$.
For one derivative, the only remaining product is
$(\partial_iv)\cdot\nabla v$. For an ordered pair of
derivatives, there are exactly three Leibniz cross-products,
each bounded using
$\|\nabla v\|_\infty$ and the full $H^2$ energy.
Bounding each by the latter energy and summing the three
first-order and nine ordered second-order identities gives
\[
\frac12\frac d{dt}\|v(t)\|_{H^2}^2
\le30\|\nabla v(t)\|_\infty\|v(t)\|_{H^2}^2.
\]
The critical Fourier estimate established in
Theorem~\ref{thm:xxxiv-gaussian} gives
\[
\int_0^\infty\|\nabla v(t)\|_\infty dt
\le C_{\rm F}\int_0^\infty\|v(t)\|_{\mathcal X^1}dt
\le 4C_{\rm F}X_0/\nu=1/100.
\]
Thus $\|v(t)\|_{H^2}\le e^{3/10}\|v_0\|_{H^2}$.
A separate exact Gaussian integration gives
\[
\|\nabla\times(\psi e_z)\|_2^2=
2(\pi/2)^{3/2},\quad
\|\nabla(\nabla\times(\psi e_z))\|_2^2=
10(\pi/2)^{3/2},\quad
\|\Delta(\nabla\times(\psi e_z))\|_2^2=
70(\pi/2)^{3/2},
\]
so the initial $H^2$ squared norm is
$(\nu/(100\pi))^2\,82(\pi/2)^{3/2}$.
The triangle and reverse triangle inequalities for
$\Delta w=\Delta v-\Delta u_B$ give the two-sided estimate.
By the source (10.21), its uniform $L^2$ bound and
Lemma~\ref{lem:xxxiii-agn},
$\|\Delta u_B(1-\tau)\|_2$ diverges at least like
$\tau^{-4A/3}$. Divide the two-sided estimate by this
diverging source norm to obtain the ratio limit.
\end{proof}
